\section{Постановка задачи и~обозначения}\label{sec:problem}
\subsection{Входы, выходы и~внешние инструменты}
Условие задачи, код, документацию и~сообщения инструментов будем представлять конечными строками. Пусть $\Sigma$ --- конечный алфавит символов, а~$\Sigma^*$ --- множество конечных строк над ним. Множество $\mathcal{X}\subseteq\Sigma^*$ содержит допустимые описания заданий. Множество $\mathcal{Y}\subseteq\Sigma^*$ содержит записи решений: текст программы или набор изменений файлов в~заранее выбранном формате. Полученную запись решения будем называть \emph{кандидатом}. Проверка может обнаружить в~этом варианте программы синтаксические ошибки.

\begin{definition}[Задача программирования]
Задача $z$ --- это кортеж $(x,\mathcal{C}_0,\mathcal{Z},\mathcal{V}_z)$, где $x\in\mathcal{X}$ --- условие, $\mathcal{C}_0\in\Sigma^*$ --- исходный контекст, $\mathcal{Z}$ --- конечное множество разрешённых инструментов, а~$\mathcal{V}_z:\mathcal{Y}\to\{0,1\}$ --- фиксированное правило внешней проверки кандидата.
\end{definition}

В~$\mathcal{C}_0$ файлы, документация и~журналы ошибок записаны с~разделителями и~именами источников. Для каждого инструмента заданы формат текстового запроса, формат ответа и~предельное время выполнения. Если время истекло, инструмент возвращает специальное сообщение. Версии инструментов и~среда фиксируются. В~этой модели проверка воспроизводима: один и~тот же кандидат в~одной среде получает одинаковый результат. Случайность внешних проверок здесь не рассматривается.

Значение $\mathcal{V}_z(y)=1$ означает, что кандидат $y$ прошёл все предусмотренные проверки за отведённое время. Значение $0$ означает отрицательный результат, ошибку или превышение времени проверки. В~журнале эти причины различаются. Успех относится к~выбранному набору проверок и~не доказывает правильность программы на всех возможных входах.

Длину строки $u$ в~токенах обозначим $\len{u}$. Для всех длин используется один зафиксированный токенизатор --- правило разбиения текста на токены. В~частности, $L=\len{\mathcal{C}_0}$ --- длина исходного контекста. Символ $\bot$ означает, что процедура завершилась без кандидата. В~этом случае она возвращает этот символ и~причину остановки. В~формальной модели это неуспешное завершение процедуры. Незавершённые экспериментальные запуски с~неизвестным исходом учитываются отдельно по правилам раздела~\ref{sec:eval}.

\subsection{Случайность, ресурсы и~качество}
Представим, что одну задачу несколько раз решают одной моделью с~теми же инструкциями и~лимитами. В~одном запуске программа может пройти все тесты с~первой попытки, в~другом модель может вернуть иной код или израсходовать больше токенов. Ответы модели, обращения к~инструментам и~их результаты составляют историю запуска. Вероятность события означает, насколько часто оно происходило бы при многократном повторении процедуры в~тех же условиях. Например, событие <<программа прошла проверку>> объединяет все истории с~таким исходом.

Для процедуры $A$ обозначим итоговую программу через $Y_A$, число потраченных токенов через $T_A$, стоимость через $C_A$, а~длительность через $\ell_A$. Расход отдельного запуска включает все обращения к~модели, в~том числе подготовку контекста и~повторы. Выполнение тестов занимает время. Если оценщик сам не обращается к~модели, её токены на тестирование не расходуются.

Средний расход на \emph{одну и~ту же задачу} обозначаем $\E[T_A]$. Его оценивают средним числом токенов по повторным запускам. Вероятность успеха оценивают долей запусков, в~которых программа прошла проверку. Если процедура штатно завершилась без программы ($\bot$), это известный неуспех. Если связь со~средой прервалась и~результат теста неизвестен, запись остаётся пропуском. Усреднение результатов по \emph{разным задачам} описано в~разделе~\ref{sec:metrics}.

Для работы задаются три лимита: $B$ токенов, $B_{\$}$ денежных единиц и~$B_{\ell}$ секунд. Перед каждым платным вызовом, в~том числе вспомогательной модели, система резервирует его максимальный разрешённый расход. Если остатка бюджета недостаточно, вызов не выполняется. Каждой операции отводится не больше оставшегося общего времени. Это правило распространяется и~на подготовку контекста, и~на внешнюю проверку. При исчерпании $B_{\ell}$ работа прекращается. Если сервис не позволяет надёжно ограничить токенный расход, $B$ остаётся целевым лимитом, а~его превышение записывается отдельно.



\section{Метод Hybrid-INoT}\label{sec:method}
\subsection{Состояние и~интроспективный проход}
В~процедуре с~внешней обратной связью одна \emph{итерация} включает интроспективный вызов и~проверку полученного кода. Выполнение трёх ролевых инструкций в~одном обращении далее называем \emph{интроспективным проходом}. Если проверка обнаруживает ошибку, следующая итерация начинает новое обращение к~модели. Перед итерацией $i$ система хранит условие задачи, подготовленный контекст, историю попыток, номер итерации и~остаток токенного бюджета. Вместе эти данные образуют состояние
\begin{formulakey}
\begin{equation}
\sigma_i=(x,\widehat{\mathcal{C}}_i,\mathcal{H}_i,i,b_i).
\end{equation}
\begin{symbols}
\sym{\sigma_i}{состояние перед вызовом}
\sym{i}{номер итерации}
\sym{x}{условие задачи}
\sym{\widehat{\mathcal C}_i}{подготовленный контекст}
\sym{\mathcal H_i}{записи предшествующих попыток, включая программы, оценки и~диагностику}
\sym{b_i}{число токенов, оставшихся в~бюджете}
\end{symbols}
\end{formulakey}
Если часть истории снова передают модели, её токены учитываются в~новом вызове.

Предельное число итераций задаётся целым $I\ge1$. За один проход модель готовит целое число кандидатов $K\ge1$. Множество текстовых планов обозначим через $\Pi\subseteq\Sigma^*$.
\begin{formulakey}
Входы и~выходы трёх действий связаны следующим образом.
\nopagebreak
\begin{align}
r_{\mathrm{plan}} &: (x,\widehat{\mathcal{C}}_i,b_i)\mapsto\pi_i\in\Pi,\\
r_{\mathrm{work}} &: (\pi_i,\widehat{\mathcal{C}}_i,K)\mapsto(y_{i,1},\ldots,y_{i,K})\in\mathcal{Y}^{K},\\
r_{\mathrm{crit}} &: ((y_{i,k})_{k=1}^{K},\widehat{\mathcal{C}}_i)\mapsto(s_{i,1},\ldots,s_{i,K})\in[0,1]^K.
\end{align}
\begin{symbols}
\sym{r_{\mathrm{plan}},\,r_{\mathrm{work}},\,r_{\mathrm{crit}}}{действия планировщика, исполнителя (worker) и~рецензента (critic)}
\sym{i,\ x,\ \widehat{\mathcal C}_i,\ b_i}{номер итерации, условие задачи, её подготовленный контекст и~оставшийся бюджет}
\sym{\pi_i\in\Pi}{текст плана из множества допустимых планов}
\sym{K,\ k}{число вариантов программы в~ответе и~номер варианта от 1 до~$K$}
\sym{y_{i,k},\ \mathcal Y^K}{текст одного варианта и~множество упорядоченных списков из~$K$ вариантов}
\sym{s_{i,k}\in[0,1]}{оценка варианта рецензентом по инструкции. Более высокое значение означает более предпочтительный вариант}
\end{symbols}
\end{formulakey}
В~запросах эксперимента роли исполнителя и~рецензента помечены словами \texttt{IMPLEMENTER} и~\texttt{REVIEWER} (приложение~\ref{app:prompt-lines}).

Все три действия задаются одной инструкцией. Если из ответа нельзя извлечь план, варианты программ и~их оценки, контроллер записывает ошибку формата.

Контроллер выбирает программу с~наибольшей оценкой рецензента. Если одинаковую максимальную оценку получили несколько программ, для проверки берётся первая из них в~порядке ответа.
\begin{formulakey}
\begin{equation*}
\begin{aligned}
 k_i^\star &= \min\{k: 1\le k\le K,\ s_{i,k}=\max_{1\le j\le K}s_{i,j}\},\\
 y_i &=y_{i,k_i^\star},\\
 s_i &=s_{i,k_i^\star},
\end{aligned}
\end{equation*}
\begin{symbols}
\sym{i,\ K}{номер итерации и~число программ в~ответе}
\sym{k,\ j}{номера вариантов в~этом ответе}
\sym{k_i^\star}{номер выбранной программы}
\sym{y_{i,k},\ s_{i,k}}{программа под номером~$k$ и~её оценка рецензентом}
\sym{y_i,\ s_i}{выбранная программа и~её оценка}
\end{symbols}
\end{formulakey}
\begin{formulakey}
После запуска программы внешняя проверка возвращает три поля.
\nopagebreak
\begin{equation*}
\begin{aligned}
 v_i &= \mathcal V_z(y_i)\in\{0,1\},\\
 M_i &\in [0,1]\cup\{\mathrm{NA}\},\\
 e_i &\in\Sigma^*,
\end{aligned}
\end{equation*}
\begin{symbols}
\sym{i,\ z,\ y_i}{номер итерации, задача и~выбранная программа}
\sym{v_i,\ \mathcal V_z}{результат проверки и~правило проверки для задачи. Значение 1 означает успех, 0 означает неуспех}
\sym{M_i}{дополнительный показатель свойств кода из раздела~
ef{sec:metrics}. Метка NA означает, что показатель недоступен}
\sym{e_i,\ \Sigma^*}{сообщение о~результате и~множество допустимых текстовых строк}
\end{symbols}
\end{formulakey}

\subsection{Выбор режима работы}\label{subsec:route-rule}
Контроллер читает требования к~организации работы до первого обращения к~модели. Обычно он объединяет планирование, написание кода и~рецензирование в~одном запросе. Если условия задачи требуют разделить исполнителей или параллельно выполнять независимые подзадачи, он выбирает отдельные обращения.
\begin{formulakey}
\begin{equation}
\mathrm{RouteMode}(z)=
\begin{cases}
\mathrm{external}, & \varphi_{\mathrm{sep}}(z)=1\ \lor\ \varphi_{\mathrm{par}}(z)=1,\\
\mathrm{internal}, & \text{в остальных случаях},
\end{cases}
\end{equation}
\begin{symbols}
\sym{z}{задача вместе с~требованиями к~её выполнению}
\sym{\varphi_{\mathrm{sep}}}{признак разделения исполнителей, например требование отдельной сессии рецензента}
\sym{\varphi_{\mathrm{par}}}{признак обязательного параллельного выполнения независимых подзадач}
\sym{\lor}{логическое <<или>>. Достаточно любого из двух требований}
\sym{\mathrm{internal},\,\mathrm{external}}{один общий вызов или отдельные обращения к~модели}
\end{symbols}
\end{formulakey}
В~пилотном сравнении использован первый признак. Если в~карточке задачи явно запрошена отдельная сессия рецензента, контроллер последовательно делает три обращения, передавая между ними план и~код. В~остальных случаях он делает один вызов. Маршрут задан до генерации, выбор не требует обращения к~LLM. Все задачи пилота допускают последовательное выполнение. Измерения приведены в~разделе~\ref{subsec:route-pilot}.

\subsection{Повторные попытки и~остановка}
Ход одного запуска и~условия остановки показаны на диаграммах приложения~\ref{app:uml}. В~конфигурации с~исправлениями результат доступного во время работы теста можно передать модели в~следующем запросе. В~общий бюджет входят все обращения к~модели. Итоговые независимые тесты запускаются после завершения генерации.

При неуспехе возвращается программа с~наибольшей оценкой рецензента среди проверенных кандидатов. Отрицательный результат её внешней проверки сохраняется. Если за все попытки получены только ошибки формата, система возвращает $\bot$.

\Needspace{8\baselineskip}
\begin{proposition}[Конечность процедуры]\label{prop:termination}
Пусть $I$ конечно. Подготовка контекста, каждый вызов модели и~каждая внешняя проверка завершаются за конечное время либо прекращаются по лимиту времени. Тогда режим с~интроспективным вызовом на рисунке~\ref{fig:uml-iteration} завершается после не более чем $I$ интроспективных проходов.
\end{proposition}
\begin{proof}
Внешний цикл содержит не более $I$ итераций. На каждой выполняется конечное число операций, ограниченных по времени. При успехе или исчерпании лимита цикл завершается раньше. Следовательно, вся процедура завершается за конечное время.
\end{proof}

Конечность процедуры не означает, что оценки рецензента сходятся или что программа верна. Из условия $s_i\in[0,1]$ не следует сближение соседних оценок. Например, последовательность $0,1,0,1,\ldots$ ограничена, но разность соседних значений по модулю всегда равна единице.

Для досрочной остановки можно использовать правило $|s_i-s_{i-1}|<\varepsilon$, где $i\ge2$, а~$\varepsilon>0$ задано заранее. Оно применяется к~двум соседним проходам, из которых удалось извлечь оценки. Это практическое правило сокращения затрат: оно может завершить поиск до того, как найдена успешная программа. В~базовой конфигурации правило выключено.

\subsection{Сжатие контекста}
Сжатие --- сокращение текста, который передаётся модели. Оно может снизить расход входных токенов, но требует собственных ресурсов и~может удалить важную деталь. В~журнал записываются длина полученного текста и~затраты на само сжатие.

Обязательные части контекста включают условие, ограничения, изменяемый код и~необходимые интерфейсы. Для остальных фрагментов можно использовать оценку релевантности, например на основе LLMLingua-2~\cite{pan2024llmlingua2}. Порядок отбора и~проверка лимита показаны на рисунке~\ref{fig:uml-context} приложения~\ref{app:uml}. В~экспериментальных сравнениях модель получала контекст без сжатия.
\FloatBarrier
